
Can we tell how a language model was trained from its benchmark scores alone?
We show that the answer is yes: a model's alignment strategy --- Direct Preference
Optimization (DPO) versus Supervised Fine-Tuning (SFT) --- can be inferred from four
standard trustworthiness benchmarks without any access to training data, at 83.3\%
leave-one-out accuracy (permutation p=0.031, n=12 models). Surprisingly, the
signal that makes this inference possible is not fairness, as alignment theory
would predict, but truthfulness: TruthfulQA MC2 carries 9.$5\times$ more discriminative
information than the next benchmark (Fisher's criterion 0.8122 vs 0.0856), with
DPO models scoring higher on truthfulness than matched SFT models
(+4.6pp per-pair mean delta) --- the opposite of the predicted direction. Fairness
benchmarks (BBQ, WinoGender) show no systematic DPO advantage (BBQ: k=3/6 pairs,
p=0.66; WinoGender: k=4/6, p=0.34). These results demonstrate the feasibility of
practical alignment auditing from public leaderboard data while challenging the
widely-held assumption that DPO preference training primarily improves fairness,
and motivate new inquiry into what preference optimization actually optimizes.

